\section{Discussion and Limitations}
\label{sec:discussion}

\method{} treats environments as persistent research assets and target trajectories as evidence for constructing them. This is a stronger modeling obligation than adapting a trusted target backend: observed traces underdetermine hidden state and unobserved actions. A program can execute correctly while training a simplified or irrelevant decision. Empty-bank initialization removes a supplied business environment, not the prior knowledge in language models, interface code, or generic examples.

Source-linked analysis is an investigation aid rather than an objective capability detector. Existing audits identify unsupported prerequisites inferred from the agent's own errors, successful traces described too favorably, and changing domain granularity across rounds. Stable task membership prevents some counting ambiguities, but composite-task scores do not isolate individual skills. Likewise, aligned goal data and bounded solvability review do not guarantee a correct verifier; author and reviewer errors can be correlated.

The present results are preliminary, single-seed and protocol-qualified. They lack fixed-library and budget-matched adaptive controls; historical training anomalies and maintenance interventions are disclosed rather than erased. Development-conditioned environment research can overfit its feedback. Although test data are excluded from the Planner, repeated operator inspection of checkpoint tests limits claims of a pristine final holdout. A final confirmatory study should freeze the method and selection rule before using a fresh test set. Most importantly, later regressions show that discovering a useful checkpoint and sustaining reliable self-improvement remain different research outcomes.
